\section{Finite Surveys}

Section 2.2 addresses how the Spherical Fourier-Bessel (SFB) decomposition framework must be modified when working with real, finite survey observations. Unlike the idealized infinite, homogeneous universe assumed in Section 2.1, actual cosmological surveys have limited volume, incomplete sky coverage, and restricted redshift ranges. This section develops the mathematical formalism to handle these observational constraints and derives key limiting cases that connect the full 3D formalism to 2D angular analysis.

\subsection{Observed Fields and Selection Functions}

In practice, astronomical surveys only observe a fraction of the total matter density field due to instrumental limitations and survey design. The relationship between the true underlying field $f(\mathbf{r})$ and the observed field $f^{\text{obs}}(\mathbf{r})$ is:

\begin{equation}
f^{\text{obs}}(\mathbf{r}) = \varphi(\mathbf{r}) f(\mathbf{r})
\label{eq:observed_field}
\end{equation}

where $\varphi(\mathbf{r})$ is the \emph{selection function} (or \emph{window function}) of the survey, which represents the survey's observational efficiency as a function of position. The selection function equals unity in regions where the survey has full sensitivity and drops to zero in unobserved regions.

The selection function $\varphi(\mathbf{r})$ can be decomposed into two components:
\begin{itemize}
\item \textbf{Radial selection function:} $\varphi_r(r)$, accounting for redshift-dependent effects such as flux limits, surface brightness limits, and survey depth variations
\item \textbf{Tangential selection function:} $\varphi_t(\theta, \phi)$, accounting for incomplete sky coverage and masking due to Galactic dust, bright stars, or other contaminants
\end{itemize}

For simplicity, this section assumes full sky coverage ($\varphi_t = 1$) and focuses on radial selection effects. The formalism naturally extends to include tangential masking via convolution with the angular 2D power spectrum.

\subsection{Selection Function Normalization}

For mathematical convenience and physical interpretation, the selection function is normalized such that:

\begin{equation}
\int d^3\mathbf{r} \, \varphi(\mathbf{r}) = V
\label{eq:selection_normalization}
\end{equation}

where $V$ is a characteristic volume of the survey defined such that $\varphi \to 1$ as $V \to \infty$. This normalization ensures that:

\begin{equation}
\lim_{V \to \infty} f^{\text{obs}}(\mathbf{r}) = f(\mathbf{r})
\label{eq:limiting_case}
\end{equation}

Physically, this says that in the limit of infinite survey volume, the observed field approaches the true underlying field, recovering the infinite-universe limit.

\subsection{Impact on Statistical Homogeneity and Isotropy}

A crucial consequence of introducing a non-uniform selection function is the breakdown of \textbf{Statistical Isotropy and Homogeneity (SIH)}, which was assumed in Section 2.1. The observed field $f^{\text{obs}}(\mathbf{r})$ is no longer homogeneous in the radial direction due to the radial selection function $\varphi_r(r)$. This means that the two-point correlation function of the observed field depends on \emph{both} the radial wavenumbers individually, rather than just their difference.

The two-point correlation function in Fourier space becomes:

\begin{equation}
\left\langle f_{\ell m}^{\text{obs}}(k) f_{\ell' m'}^{\text{obs}*}(k') \right\rangle = C_\ell^{\text{obs}}(k, k') \delta_{\ell \ell'} \delta_{m m'}
\label{eq:observed_corr_func}
\end{equation}

where the observed SFB power spectrum $C_\ell^{\text{obs}}(k, k')$ now depends explicitly on both $k$ and $k'$, as opposed to only on the difference $|k - k'|$ as in the infinite-universe case (Equation 5 in Section 2.1).

\subsection{Observed SFB Power Spectrum}

The observed SFB power spectrum can be expressed in terms of the true 3D matter power spectrum $P(k)$ via a convolution with a window function:

\begin{equation}
C_\ell^{\text{obs}}(k, k') = \frac{2}{\pi} \int_0^\infty dk'' \, k''^2 \, P(k'') \, W_\ell(k, k'') \, W_\ell(k', k'')
\label{eq:obs_power_spectrum}
\end{equation}

This formula shows that the observed power spectrum is a smoothed version of the true power spectrum, with the smoothing kernel determined by the survey's window function. The true power spectrum $P(k)$ is thus recovered (up to quadratic factors and normalization) from the observed SFB power spectrum through deconvolution.

\subsection{Window Function Definition}

The \emph{window function} $W_\ell(k, k')$ encodes the impact of the survey's spatial coverage on the measured power spectrum. It is defined as:

\begin{equation}
W_\ell(k, k') = \int_0^\infty dr \, r^2 \, \varphi(r) \, j_\ell(kr) \, j_\ell(k'r)
\label{eq:window_function}
\end{equation}

This integral represents the overlap of two SFB basis functions $j_\ell(kr)$ and $j_\ell(k'r)$, weighted by the selection function. Physically, the window function characterizes how the survey geometry couples modes with different wavenumbers $k$ and $k'$.

Key properties of the window function:
\begin{itemize}
\item \textbf{Symmetry:} $W_\ell(k, k') = W_\ell(k', k)$ (hermitian)
\item \textbf{Diagonal dominance:} The window function is largest along the diagonal $k = k'$ and falls off rapidly for $|k - k'| \gg W_\ell^{-1}$
\item \textbf{Wavenumber widths:} The characteristic width in wavenumber space is $\Delta k \sim 1/r_{\text{max}}$, where $r_{\text{max}}$ is the radial extent of the survey
\item \textbf{Normalization:} The window function is normalized such that $\int dk'' \, W_\ell(k, k'') \sim 1$ for a uniform selection function
\end{itemize}

\subsection{Diagonal Approximation}

In practice, the observed power spectrum $C_\ell^{\text{obs}}(k, k')$ falls off rapidly away from the diagonal $k = k'$. Therefore, it is often convenient to work with the \emph{diagonal approximation}:

\begin{equation}
\bar{C}_\ell^{\text{obs}}(k) \equiv C_\ell^{\text{obs}}(k, k)
\label{eq:diagonal_approx}
\end{equation}

This diagonal power spectrum is easier to interpret and compute, though it discards information about mode coupling. The diagonal power spectrum is related to the true power spectrum by:

\begin{equation}
\bar{C}_\ell^{\text{obs}}(k) = \frac{2}{\pi} \int_0^\infty dk'' \, k''^2 \, P(k'') \left[ W_\ell(k, k'') \right]^2
\label{eq:diagonal_relation}
\end{equation}

\subsection{Three Special Cases for Selection Functions}

The paper then analyzes three idealized selection function cases that provide analytical insight and serve as benchmarks for numerical studies.

\subsubsection{Gaussian Selection Function}

The first canonical case is a \textbf{Gaussian radial selection function}:

\begin{equation}
\varphi(r) = \exp\left( -\left( \frac{r}{r_0} \right)^2 \right)
\label{eq:gaussian_selection}
\end{equation}

where $r_0$ is a characteristic radius parameter. The normalization constant is $V = \pi^{3/2} r_0^3$. Physically, this represents a survey with smooth, centrally-concentrated sensitivity that falls off exponentially at large distances. This might approximate surveys with flux limits that preferentially detect nearby objects.

For the Gaussian selection function, the window function can be evaluated analytically:

\begin{equation}
W_\ell(k, k') = \frac{\pi r_0^2}{k} \sqrt{\frac{\pi r_0^2 k}{4}} \exp\left( -r_0^2 \frac{k^2 + k'^2}{4} \right) I_{\ell + 1/2}\left( \frac{r_0^2 kk'}{4} \right)
\label{eq:gaussian_window}
\end{equation}

where $I_\nu(x)$ is the \emph{modified Bessel function of the first kind} of order $\nu$. The modified Bessel function is defined by the integral:

\begin{equation}
I_\nu(x) = \frac{1}{\pi} \int_0^\pi d\theta \, e^{x \cos\theta} \cos(\nu\theta)
\label{eq:modified_bessel_def}
\end{equation}

The modified Bessel function $I_\nu(x)$ is distinct from the standard Bessel function $J_\nu(x)$ and the spherical Bessel functions $j_\ell(x)$ discussed in Section 2.1. For small arguments $x \ll 1$, we have $I_\nu(x) \approx (x/2)^\nu / \Gamma(\nu+1)$, while for large arguments $x \gg 1$, the asymptotic form is $I_\nu(x) \approx e^x / \sqrt{2\pi x}$.

The analytical form of the Gaussian window function (Equation \ref{eq:gaussian_window}) greatly simplifies calculations. Rather than evaluating the convolution integral (Equation \ref{eq:obs_power_spectrum}) numerically for each pair of wavenumbers, one can use the analytical expression directly. This efficiency gain is crucial for computational feasibility in data analysis pipelines.

The paper notes (with reference to Appendix A) that for large arguments of the modified Bessel function, special numerical techniques are required to avoid overflow or underflow errors. However, these issues can be mitigated using recurrence relations and exponentially-scaled versions of the Bessel function.

\subsubsection{Radial Limit: Full Radial Coverage}

The second limiting case corresponds to a survey with \textbf{complete radial coverage} across all distances:

\begin{equation}
\varphi(r) = 1 \quad \forall r
\label{eq:full_radial_coverage}
\end{equation}

In this limit, the window function becomes:

\begin{equation}
W_\ell(k, k') = \frac{\pi}{2k} \delta(k - k')
\label{eq:radial_limit_window}
\end{equation}

This is a \emph{delta function} in wavenumber space, meaning there is no coupling between different wavenumber modes. Substituting this into Equation \ref{eq:obs_power_spectrum}:

\begin{equation}
C_\ell^{\text{obs}}(k, k') = C_\ell(k) \delta(k - k') = P(k) \delta(k - k')
\label{eq:radial_limit_power}
\end{equation}

we recover the infinite-universe result from Section 2.1: the SFB power spectrum depends only on the radial wavenumber $k$. This limit represents the \emph{radial domain}: a regime where the 3D power spectrum is fully determined by radial information and independent of angular information.

The radial limit can also be derived from the Gaussian selection function (Equation \ref{eq:gaussian_selection}) by taking the formal limit $r_0 \to \infty$. In practice, this limit is achieved when the condition:

\begin{equation}
r_0 \, k \gg \sqrt{2\ell(\ell+1)}
\label{eq:radialisation_condition}
\end{equation}

is satisfied. This is the fundamental condition for \emph{radialisation}: the survey must extend over distances much larger than the characteristic scale $\sqrt{2\ell(\ell+1)} / k$ associated with the mode. For modes with large wavenumber $k$ or low angular momentum $\ell$, radialisation occurs at smaller survey volumes.

\subsubsection{Tangential Limit: Thin-Shell Selection}

The third limiting case corresponds to a \textbf{thin-shell radial selection function}, where the survey observes only a narrow shell of thickness $\Delta r \to 0$ at a fixed distance $r_*$:

\begin{equation}
\varphi(r) = r_* \, \delta(r - r_*)
\label{eq:thin_shell_selection}
\end{equation}

The delta function captures the fact that observations are restricted to a single radial distance. The normalization is chosen such that $V = 4\pi r_*^3$ (corresponding to a spherical shell). In this limit, the observed SFB power spectrum becomes:

\begin{equation}
C_\ell^{\text{obs}}(k, k') = \frac{2}{\pi} k k' \, j_\ell(k r_*) j_\ell(k' r_*) \int_0^\infty dk'' \, k''^2 \, P(k'') \, r_*^3 \, j_\ell(k'' r_*)
\label{eq:thin_shell_power}
\end{equation}

This can be rewritten more compactly by factoring out the $k$-dependent terms. The key insight is that this 3D thin-shell result is intimately connected to 2D angular analysis, the subject of the next development.

\subsection{Connection to 2D Projected Fields and Spherical Harmonic Analysis}

Real surveys observe galaxies projected onto the celestial sphere with redshift information. This is fundamentally a 2D observable in angular coordinates $(\theta, \phi)$ combined with radial information from redshift. To bridge the 3D SFB formalism with conventional 2D harmonic analysis, we consider the \emph{2D projected field}:

\begin{equation}
f^{2D}(\theta, \phi) = \int_0^\infty dr \, r^2 \, \varphi(r) \, f(\mathbf{r})
\label{eq:projected_2d_field}
\end{equation}

This integral projects the 3D field onto the sky by summing over all radial distances, weighted by the selection function. Expanding this 2D field in spherical harmonics:

\begin{equation}
f^{2D}(\theta, \phi) = \sum_{\ell=0}^\infty \sum_{m=-\ell}^\ell f_{\ell m}^{2D} Y_{\ell m}(\theta, \phi)
\label{eq:sh_expansion_2d}
\end{equation}

where the spherical harmonic coefficients are:

\begin{equation}
f_{\ell m}^{2D} = \int d\Omega \, Y_{\ell m}^*(\theta, \phi) \, f^{2D}(\theta, \phi)
\label{eq:sh_coeff_2d}
\end{equation}

The 2D power spectrum (angular power spectrum) is the two-point statistics:

\begin{equation}
\left\langle f_{\ell m}^{2D} f_{\ell' m'}^{2D*} \right\rangle = C_\ell^{2D} \delta_{\ell\ell'} \delta_{m m'}
\label{eq:2d_power_spectrum}
\end{equation}

which depends only on the multipole index $\ell$ due to isotropy on the sphere.

\subsection{2D Window Function and Relation to 3D Power Spectrum}

The 2D power spectrum is related to the 3D matter power spectrum through its own window function:

\begin{equation}
C_\ell^{2D} = \frac{2}{\pi} \int_0^\infty dk \, k^2 \, P(k) \, \left[ W_\ell^{2D}(k) \right]^2
\label{eq:2d_cl_relation}
\end{equation}

where the \emph{2D window function} is defined as:

\begin{equation}
W_\ell^{2D}(k) = \int_0^\infty dr \, r^2 \, \varphi(r) \, j_\ell(kr)
\label{eq:2d_window_function}
\end{equation}

Note that the 2D window function is a one-dimensional object (depending only on the single wavenumber $k$), in contrast to the 3D window function $W_\ell(k, k')$ which depends on two wavenumbers.

Due to the normalization choice in Equation \ref{eq:selection_normalization}, the 2D window function $W_\ell^{2D}(k)$ has units of volume (length$^3$). Physically, $W_\ell^{2D}(k)$ represents the radially-integrated sensitivity of the survey for observing modes with wavenumber $k$ and angular momentum $\ell$.

\subsection{Thin-Shell Case: Connecting 3D and 2D Analysis}

For the thin-shell selection function (Equation \ref{eq:thin_shell_selection}), the 2D window function simplifies to:

\begin{equation}
W_\ell^{2D}(k) = r_*^3 \, j_\ell(k r_*)
\label{eq:thin_shell_2d_window}
\end{equation}

This is simply the spherical Bessel function evaluated at the fixed shell radius $r_*$, scaled by the shell volume $r_*^3$.

For the same thin-shell selection, the 2D power spectrum is:

\begin{equation}
C_\ell^{2D} = \frac{2}{\pi} \int_0^\infty dk \, k^2 \, P(k) \, r_*^6 \, \left[ j_\ell(k r_*) \right]^2
\label{eq:thin_shell_2d_cl}
\end{equation}

Comparing with the 3D thin-shell power spectrum (Equation \ref{eq:thin_shell_power}), we can establish a direct relationship:

\begin{equation}
C_\ell^{\text{obs}}(k, k') = \frac{2}{\pi} k k' \, j_\ell(k r_*) j_\ell(k' r_*) \, C_\ell^{2D}
\label{eq:3d_2d_relation}
\end{equation}

The factor $\frac{2}{\pi} k k' j_\ell(k r_*) j_\ell(k' r_*) C_\ell^{2D}$ is a \emph{geometric factor} that does not depend on the underlying matter power spectrum $P(k)$ or the field $f$. Therefore, all the cosmological information contained in the 3D SFB power spectrum for a thin-shell survey is captured by the 2D angular power spectrum $C_\ell^{2D}$.

This is a profound result: for surveys observing a thin redshift slice, the 3D structure cannot be resolved, and the full 3D information collapses to 2D information. The 3D SFB coefficients become products of spherical Bessel functions evaluated at the shell radius, times the 2D power spectrum. This establishes the formal connection between the full 3D spherical analysis (the focus of this paper) and conventional 2D angular harmonic analysis used in CMB and weak lensing studies.

\subsection{Physical Interpretation and Limitations}

The three cases analyzed above represent limiting regimes:

\begin{itemize}
\item \textbf{Radial limit:} Surveys with deep radial coverage allow measurements of the 3D power spectrum as a function of $k$. No angular resolution in 3D is possible.
\item \textbf{Tangential limit:} Surveys observing thin redshift shells yield 2D angular power spectra. No radial 3D structure can be resolved.
\item \textbf{Intermediate case (Gaussian):} Real surveys occupy an intermediate regime where both radial and angular information contribute. The window function $W_\ell(k, k')$ characterizes mode-coupling effects due to the finite survey volume and incomplete coverage.
\end{itemize}

For intermediate cases, the mode-coupling encoded in the off-diagonal elements of $W_\ell(k, k')$ becomes important. Recovering the true 3D power spectrum $P(k)$ requires accounting for these mode-coupling corrections, which can be computationally expensive. Various techniques exist to address this: direct inversion using eigenmodes of the window matrix, optimal quadratic estimators, or Bayesian inference methods that marginalize over the unknown power spectrum.

\subsection{Implications for BAO Analysis}

Baryon Acoustic Oscillations (BAOs) are subtle oscillatory features in the matter power spectrum at wavenumbers $k \sim 0.1 \, h \, \text{Mpc}^{-1}$. The selection function and window function effects have profound impacts on BAO measurements:

\begin{itemize}
\item \textbf{Suppression of high-$k$ modes:} Finite survey volume limits the maximum wavenumber $k_{\max} \sim 2\pi / L_{\min}$, where $L_{\min}$ is the smallest linear scale sampled.
\item \textbf{Coupling of BAO scales:} The window function couples nearby wavenumber modes, smearing the BAO feature. The characteristic smearing scale is $\Delta k \sim \pi / L$, where $L$ is the survey length.
\item \textbf{Radialisation vs. tomography:} In the radial limit, BAO measurements are restricted to radial BAO (acoustic horizon). In the 2D limit, angular BAO can be measured. Intermediate surveys access both.
\item \textbf{Reconstruction gains:} BAO reconstruction methods can partially reverse mode-coupling effects, recovering some power in the BAO feature that was redistributed by the window function.
\end{itemize}

These effects are central to the motivation for the spherical Fourier-Bessel method: by explicitly accounting for the 3D window function, one can extract more cosmological information from finite surveys than would be available in traditional 2D approaches. The 3D formalism makes the mode-coupling explicit and amenable to treatment via the window function deconvolution techniques discussed throughout the paper.

